\subsection{Central banks, the policy path and the curve}

\begin{qfdefinition}[Central-bank policy framework: mandates, conventional tools, the balance sheet and money aggregates, the zero lower bound, quantitative easing and its measured effects, crisis lending facilities, and hawks and doves]\label{def:central-bank-policy-framework}
\qffield{Statement} Gliner's chapter on central banks \citep[chap.~11, PDF pp.~355--400]{s7-gliner}:
\begin{itemize}
\item \emph{Goals} (sec.\ ``Monetary Policy Goals'', pp.~356--357): price stability above all, with the Fed's statutory goals of maximum employment, stable prices and moderate long-term rates, the Bank of England's monetary and financial stability, and the SNB's price stability under its law.
\item \emph{Tools} (sec.\ ``Tools Used by Central Banks'', pp.~358--361): the policy rate, the overnight rate at which banks lend reserves, which he calls ``the single strongest action''; the reserve requirement, with a money multiplier of $1/r$ (10 at 10 percent); open market operations and repos; discount-window lending; and currency intervention, \emph{sterilised} when the central bank offsets the effect on the monetary base and \emph{unsterilised} when it does not (pp.~360--361).
\item \emph{Monetary base and money supply} (pp.~365--367): the base is currency plus reserves; the quantity identity $MV=PY$; aggregates M1--M3 in Table~11.1, with M2 ``historically'' a predictor of nominal income.
\item \emph{Zero lower bound and liquidity trap} (p.~370), and \emph{quantitative easing} (pp.~371--374): purchases of government bonds and MBS financed with reserves; Table~11.4, from Bauer's San Francisco Fed letter of May 2012, gives the cumulative change in the 10-year Treasury yield on announcement days of $-1.00$ percentage point for QE1 (five dates, 25 November 2008 to 18 March 2009), $-0.14$ for QE2 (three dates in 2010) and $-0.08$ for the maturity extension programme (21 September 2011) (p.~373).
\item \emph{Communication} (pp.~375--379): hawks and doves, with a dove's hawkish remark moving markets more than a hawk's, and the Bank of England's minutes as the example (\hyperref[meth:macro-event-queries]{event queries}).
\item \emph{Fed crisis programmes} (appendix, pp.~396--398): central-bank liquidity swap lines (first announced 12 December 2007); the Term Auction Facility (28-day, later 84-day, loans auctioned to avoid the discount-window stigma); the Term Securities Lending Facility (Treasuries lent to primary dealers for 28 days against other collateral); the Primary Dealer Credit Facility; the AMLF; the Commercial Paper Funding Facility; the Term Asset-Backed Securities Loan Facility; and agency MBS purchases of \$1.25 trillion in QE1 (p.~371).
\end{itemize}

\qffield{Notation} $r$ required reserve ratio; $M$ money, $V$ velocity, $P$ price level, $Y$ real output; M0 to M3 as each central bank defines them.

\qffield{Conditions} \emph{Announcement-day effects are not total effects.} Bauer's letter (FRBSF Economic Letter 2012-16, 21 May 2012) gives the same figures, $-1.00$, $-0.14$ and $-0.08$ for the 10-year Treasury yield, as the ``Cumulative change of key interest rates in percentage points over the announcement days'' of each programme, and warns that the event study ``misses effects that occur on days other than those when key announcements are made'' and that ``other events may occur on announcement days'' \citep[Table~1 and text]{bauer-2012}; Gliner reproduces the table faithfully. These figures sum one-day yield changes on announcement dates; they measure the surprise component on those days, not the cumulative effect of the purchases, and depend on which dates are chosen.

\emph{The underlying event study}. Gagnon, Raskin, Remache and Sack's working-paper version sums eight baseline announcements from November 2008 to November 2009: the 10-year Treasury yield, 10-year agency debt yield and current-coupon agency MBS yield fell 91, 156 and 113~bp; on all other days since November 2008 the 10-year Treasury yield and the swap rate ``increased nearly 100 basis points'', leaving them ``roughly unchanged over the entire period'', which the authors attribute to other factors such as the fiscal outlook \citep[pp.~17--18 (PDF pp.~19--20)]{nyfed-sr441}. Their time-series estimate is that the purchases ``reduced the 10-year term premium by between 50 and 100 basis points'' (p.~21), and their summary range is ``somewhere between 30 and 100 basis points'' (p.~28). So the event-day sums, the net change in yields and the model-based term-premium effect are three different numbers.

\emph{Money multiplier.} $1/r$ is a textbook upper bound; with interest on reserves and ample reserves, bank lending does not follow reserves mechanically. \emph{Definitions change.} Gliner's U.S.\ M1--M3 are as of 2013; current H.6 definitions differ.

\qffield{Practice} Three policy regimes (conventional, with the rate as the tool; unconventional, at the lower bound, working through asset prices, expectations and FX; crisis, where liquidity provision dominates), four QE channels (portfolio rebalancing, term-premium and credit-spread compression, signalling, exchange rate), and mandates that include currency stability for the SNB and PBOC are a reasonable synthesis, but an interpretive framing, not Gliner's text.

\qffield{Pitfalls} Open-market operations are not the only mechanism affecting the monetary base; lending facilities, FX intervention and interest on reserves also change it. That at the ZLB yield curves flatten naturally, carry trades distort, volatility is suppressed, forward guidance becomes critical and fiscal policy matters more are untested assertions. Reading Bauer's smaller later numbers as declining effectiveness over successive rounds is not the only reading; smaller surprises would give the same pattern.

\qffield{See also} \hyperref[mod:taylor-rule]{Taylor rule}; \hyperref[rule:major-central-banks-current]{major central banks}; \hyperref[meth:policy-path-ois-futures]{policy path from OIS and futures}; \hyperref[def:currency-regimes-intervention]{currency regimes and intervention}; \hyperref[def:money-market-reference-rates]{money-market reference rates}.

\qffield{Further reading} Large-scale asset purchases \citep{rt-gagnon-lsap-2011}; its March 2010 working-paper version \citep{nyfed-sr441} provides the event-study estimates. The published version may differ from the working paper.

\end{qfdefinition}

\begin{qfmodel}[The Taylor rule: Taylor's 1993 equation, its 2 percent equilibrium real rate and coefficients, the Taylor principle, and Gliner's printed variants]\label{mod:taylor-rule}
\qffield{Statement} Taylor's rule for the federal funds rate is \citep[p.~202, eq.~(1)]{taylor-1993}
\[ r=p+0.5\,y+0.5\,(p-2)+2, \]
where $r$ is the federal funds rate, $p$ the rate of inflation over the previous four quarters and $y$ the percent deviation of real GDP from target; the equilibrium real rate and the inflation target are both 2 percent. The general form used is $i_t=r^*+\pi_t+\alpha(\pi_t-\pi^*)+\beta(y_t-y^*_t)$, with $r^*=2$, $\pi^*=2$ and $\alpha=\beta=0.5$ in Taylor's version. Gliner quotes the Kansas City Fed's verbal definition and prints the formula as an image \citep[chap.~11, sec.\ ``Taylor Rule'', PDF pp.~361--362]{s7-gliner}.

\qffield{Notation} $i_t$ (Taylor's $r$) nominal policy rate; $\pi_t$ inflation; $\pi^*$ target; $y_t-y^*_t$ output gap in percent; $r^*$ equilibrium real rate.

\qffield{Conditions} The response of the nominal rate to inflation is $1+\alpha=1.5$, greater than one: the \emph{Taylor principle}, under which a rise in inflation raises the real rate. On target with no gap the rule gives 4 percent; inflation of 4 percent gives 7 percent; a $-2$ percent gap gives 3 percent. Taylor fitted the rule to U.S.\ policy over 1987--1992 as a description and benchmark, not as a law central banks follow.

\qffield{Practice} Gliner suggests forecasting the rule's inputs with the market's one-year CPI forecast and an output measure, and comparing the result with fed funds futures (p.~362).

\qffield{Pitfalls} The neutral-rate estimate $r^*$ is model-dependent; Taylor used 2 percent, and current estimates vary by model. One of Gliner's two printed forms reads ``$p+1/(2y)$'', a typesetting slip for $\tfrac12 y$: at $y=2$ it gives 4.25 instead of 5.00 (checks, section~5). Inputs measured with revised data overstate how well the rule fits in real time.

\qffield{See also} \hyperref[def:central-bank-policy-framework]{central-bank policy}; \hyperref[meth:policy-path-ois-futures]{policy path from OIS and futures}.

\qffield{Further reading} Taylor (1993).

\end{qfmodel}

\begin{qfrule}[Major central banks as of 16 October 2026: Fed, ECB, Bank of England, Bank of Japan and SNB mandates, committees and meeting frequency; ECB crisis tools and TARGET2; communication and the Bank of England's 2012 minutes]\label{rule:major-central-banks-current}
\qffield{Statement} \emph{Federal Reserve.} The FOMC has twelve voting members, the seven members of the Board of Governors, the New York Fed's president and four of the other eleven Reserve Bank presidents in one-year rotation, and meets eight times a year \citep[page updated 7 October 2026]{fomc-page}. Its inflation goal is 2 percent on the annual change in the PCE price index \citep[pp.~1--2]{fomc-longer-run-goals}.

\emph{ECB.} The Governing Council comprises the six members of the Executive Board and the governors of the 21 euro-area national central banks, and takes monetary policy decisions every six weeks \citep[page dated 10 September 2026]{ecb-governing-council}; its strategy targets 2 percent HICP inflation over the medium term, symmetrically \citep[web page]{ecb-strategy}. \emph{Bank of England.} The Monetary Policy Committee has nine members, four of them external \citep[web page]{boe-mpc-page}; ``The Monetary Policy summary and minutes are published at 12 noon on a Thursday eight times a year'' \citep[page last updated 31 January 2023]{boe-mps-minutes}.

\emph{Bank of Japan.} The Policy Board has nine members, the Governor, two Deputy Governors and six other members \citep[web page]{boj-policy-board}. \emph{SNB.} Price stability means a rise in the consumer price index of less than 2 percent a year, assessed quarterly \citep[web page]{snb-strategy}.

\qffield{Conditions} \emph{Historical, as Gliner gives them} \citep[chap.~11, sec.\ ``Central Banks'', PDF pp.~380--393]{s7-gliner}. Four to five of the eight FOMC meetings had projections and a press conference (p.~380). In 2025 every meeting had a press conference, and in 2026 the four meetings marked for a Summary of Economic Projections are those of March, June, September and December, with a press conference after each meeting held so far \citep[2025 and 2026 calendars]{fomc-calendars}. The ECB's Governing Council had 23 members and ``usually meets twice a month'' (p.~383).

The three-year LTROs allotted \texteuro529.5 billion in December 2011 and \texteuro489.2 billion on 29 February 2012 (p.~384). OMT was announced on 6 September 2012, with the SMP terminated, for one- to three-year bonds of countries under strict conditionality, with ``no limits to size'' and full sterilisation (pp.~386--387). German credit exposure through TARGET2 to countries like Spain had risen ``to around \texteuro700 billion'' (pp.~387--388). The SNB introduced a minimum EUR/CHF rate of 1.20 on 6 September 2011 (p.~391).

The Bank of England raised its asset purchases by \pounds50 billion to \pounds375 billion on 5 July 2012 (Table~11.3, p.~372), after June minutes showing a 5--4 vote against more QE, against 8--1 in May (pp.~377--379).

\qffield{Practice} Membership changes with euro-area enlargement: Gliner's 23 (six plus seventeen governors) and 25 (six plus nineteen) were each right for their time; 27 is current. Communication channels, list (statements after each meeting, minutes, press conferences, projections such as the Fed's dot plot, speeches), are standard; their current formats should be read from each bank's site.

\qffield{Pitfalls} \emph{Historical conventions.} ECB ``HICP $<2\%$'' is an earlier formulation (Gliner: the ECB ``aims at inflation rates of below, but close to, 2'' percent, p.~356); the current page states a 2 percent target that is ``symmetric'' \citep[web page]{ecb-strategy}. \emph{Minutes}. ``Released with lag (2 to 3 weeks)'' fits the Bank of England only before August 2015: the Warsh review of December 2014 described minutes released ``about a fortnight after each meeting'' \citep[PDF p.~10]{warsh-2014}; the July 2015 decision was published on 9 July and its minutes on 22 July, and from August 2015 the Monetary Policy Summary and minutes were published together, on 6 August 2015 \citep[2015 documentation page]{boe-mpc-docs-2015}. FOMC minutes follow about three weeks after each meeting (those of 15--16 September 2026 on 7 October 2026) \citep[2026 calendar]{fomc-calendars}.

\emph{TARGET2.} Gliner calls TARGET2 a platform ``for settling almost all heavily traded securities'' (p.~387); it is the Eurosystem's large-value payment system, and securities settlement is the separate T2S service. \emph{Later central-bank history}. The Bank of Japan ended its framework of Quantitative and Qualitative Monetary Easing with Yield Curve Control and the negative interest rate policy on 19 March 2024, guiding the overnight call rate to around 0 to 0.1 percent and discontinuing ETF and J-REIT purchases, by a 7--2 vote on the rate guideline \citep[statement of 19 March 2024, pp.~1--2]{boj-2024-framework}; its yield curve control and negative rates are therefore history, and the claim that it pioneered QE is unverified. The SNB discontinued the EUR/CHF minimum exchange rate and lowered its rate to $-0.75$ percent on 15 January 2015 \citep[press release title and date]{snb-2015-floor}.

The ECB's asset purchase programme began purchases in October 2014, with net purchases of \texteuro60 billion a month from March 2015, and reinvestments were discontinued from July 2023 \citep[APP page]{ecb-app}; on 21 July 2022 the ECB said OMT ``is part of the Eurosystem's toolkit'' and that the Governing Council ``retains discretion to conduct'' it \citep[Transmission Protection Instrument release]{ecb-tpi-2022}. \emph{OMT's use}. In \emph{Gauweiler} the Court of Justice recorded that ``more than two years after the programme \dots\ was announced, that programme has not been implemented, the Governing Council taking the view that its activation was not justified by the economic situation of the euro area'' \citep[Case C-62/14, judgment of 16 June 2015, para.~84, via CELLAR]{cjeu-gauweiler-2015}, and its press release says the programme ``has still not been activated'' \citep[Press Release No~70/15]{cjeu-gauweiler-2015}.

That establishes non-activation to June 2015; for later years the support is the ECB's 2022 description of OMT as still ``part of the Eurosystem's toolkit'' (above) and the absence of any activation in the sources, so ``never actually used'' is well supported but rests partly on an absence. \emph{TLTROs and the renminbi band}. The ECB announced three series of targeted longer-term refinancing operations, on 5 June 2014, 10 March 2016 (TLTRO~II, ``a series of four \dots\ each with a maturity of four years'') and 7 March 2019 (TLTRO~III, ``a series of ten \dots\ each with a maturity of three years''), the third with voluntary early repayment, after one year from September 2021 under Decision ECB/2020/13 \citep[TLTRO page]{ecb-tltro}.

The PBOC's Public Announcement [2014] No.~5 widened the interbank band of the renminbi against the U.S.\ dollar: it ``is enlarged from 1 percent to 2 percent'', ``Effective from 17 March 2014'' \citep[English text on the State Council's site, attributed to \texttt{pbc.gov.cn}]{pboc-2014-5}. Whether a later notice changed the band was not established. The PBOC's English regime page states only that ``China pursues a managed floating exchange rate regime based on market supply and demand with reference to a basket of currencies'', with no band width \citep[Exchange Rate Regime page]{pboc-en-2026}.

\emph{Reserve requirements}. The PBOC's required-reserves page lists cuts announced on 9 July 2021, 6 December 2021, 15 April 2022, 17 March 2023, 24 January 2024, 27 September 2024 and 7 May 2025, among others; its English pages give no complete history. The 2024 cut lowered the required reserve ratio by 0.5 percentage point, effective 27 September 2024, ``excluding those that are already subject to an RRR of 5 percent'', leaving a weighted average of ``about 6.6 percent''. The 2025 cut, also 0.5 point, took effect on 15 May 2025, with a further 5 points for auto-finance and financial-leasing companies \citep[Required Reserves page and the two announcements]{pboc-en-2026}.

\emph{The August 2015 fixing change} (the PBC statement as reproduced in NAFMII's weekly update): ``Effective from 11 August 2015, the quotes of central parity that market makers report to the CFETS daily before market opens should refer to the closing rate of the inter-bank foreign exchange market on the previous day, in conjunction with demand and supply condition in the foreign exchange market and exchange rate movement of the major currencies'' \citep[p.~1]{nafmii-2015-08-14}. The band's current width and the reserve-requirement history are unverified.

\emph{Inflation forecasts.} The Bank of England published its first \emph{Inflation Report} in February 1993, with a two-year inflation projection and a shaded uncertainty band built from the previous ten years' forecast errors, and its first fan chart in February 1996 \citep[pp.~30--31]{boe-qb-1998-fan-chart}. That dates the Bank's own practice; it does not establish that the Bank was the first major central bank to publish inflation forecasts, which needs other central banks' dates (the Reserve Bank of New Zealand, for example).

\qffield{See also} \hyperref[def:central-bank-policy-framework]{central-bank policy}; \hyperref[data:us-macro-release-calendar]{release calendar}; \hyperref[def:currency-regimes-intervention]{currency regimes} (the SNB floor).

\qffield{Further reading} Each bank's own site for current decisions; Gliner, chap.~11, for the 2008--2013 record.

\end{qfrule}

\begin{qfmethod}[Reading the policy path: fed funds futures, overnight index swaps, meeting-dated OIS probabilities, and Eurodollar futures as a historical LIBOR contract converted to SOFR in 2023]\label{meth:policy-path-ois-futures}
\qffield{Statement} \emph{Fed funds futures} settle on 100 minus the average daily effective fed funds rate for the contract month, on a \$5 million face value \citep[chap.~9, sec.\ ``Federal (Fed) Funds'', PDF pp.~257--259, Table~9.1]{s7-gliner}. An \emph{overnight index swap} exchanges a fixed rate for the compounded overnight rate (the effective fed funds rate in the United States; in Gliner's time EONIA, SONIA and TOIS elsewhere); OIS are ``great tools to hedge or speculate in central bank action'', usually under two years, with customisable dates (sec.\ ``Overnight Indexed Swaps (OIS)'', pp.~260--261). Gliner's example: with Bank Rate at 0.50 percent, a December MPC-dated OIS at 0.625 prices ``a 50 percent chance for a hike'' of 25~bp ($(0.625-0.50)/0.25$) (p.~261).

\emph{Eurodollar futures}, as Gliner describes them in 2013--2014 (sec.\ ``London Interbank Offered Rate (LIBOR)'', pp.~254--255): \$1,000,000 notional on three-month LIBOR, price 100 minus the rate, one basis point worth \$25, March, June, September and December months, and colour codes by year (white, red, green, blue, gold and beyond).

\qffield{Notation} $P=100-R$; for a 90-day period, $\$1{,}000{,}000\times0.0001\times90/360=\$25$ per basis point.

\qffield{Conditions} \emph{Probabilities from OIS.} The $(\text{OIS}-\text{current})/\text{step}$ reading holds when the swap's period starts at the meeting's effective date (a meeting-dated OIS) and only one move of one size is possible. If part of the period precedes the meeting, the average is diluted: with 60 percent of the period after the meeting, the same quote implies 83 percent (checks, section~6). Term premia, the spread between the overnight index and the policy rate, and the possibility of a larger move make the reading approximate.

\qffield{Practice} \emph{The end of LIBOR and of Eurodollar futures.} On 5 March 2021 ICE Benchmark Administration announced that it could not publish the main U.S.\ dollar LIBOR settings on a representative basis beyond 30 June 2023, and the FCA confirmed that three-month U.S.\ dollar LIBOR would cease to be representative immediately after that date; SOFR, selected by the ARRC in 2017 and published since April 2018, became the successor, with CME SOFR futures from May 2018; CME converted open Eurodollar futures and options into three-month SOFR futures and options in April 2023 (the CBOT filing gives ``April 23, 2023'', a date corrected below); synthetic one-, three- and six-month U.S.\ dollar LIBOR was published until 30 September 2024 and barred from new contracts and cleared derivatives \citep[secs.~1--2, pp.~1--5]{cbot-23-214}.

The ARRC records that new use of U.S.\ dollar LIBOR had stopped after 31 December 2021 \citep[web page]{arrc-sofr-transition}; the end of the panel and \texteuro{}STR's replacement of EONIA are in \hyperref[def:money-market-reference-rates]{money-market reference rates}.

\emph{The conversion date}. CME Clearing Advisory 23-082 (15 March 2023) states that ``following close of business on Friday, April 14, 2023, CME will convert specific Eurodollar Futures and Options open interest to corresponding SOFR positions, at a ratio of 1:1, for Trade Date Monday, April 17, 2023'', with the transformation run at about 10~a.m.\ CST on Sunday 16 April; it covered Eurodollar futures and options from the September 2023 expiry onwards and the one- to five-year mid-curve options (codes E0 and E2--E5, June 2023 the last listed), while the May and June 2023 contracts stayed listed to expiry (the June futures settling on 19 June 2023 against that day's LIBOR) \citep[p.~1]{cme-23-082}.

The ARRC's minutes of 9 March 2023 give the same 14 April \citep[p.~1]{arrc-minutes-2023-03-09}, as did the CFTC chairman's keynote in advance (``CME intends to convert all eligible Eurodollar futures and options to SOFR equivalent contracts on April 14, 2023'') \citep[keynote text]{cftc-behnam-libor}. A later CME filing reports that ``In April 2023, CME successfully converted 7.5 million contracts of Eurodollar futures and options open interest'', citing a ``CME Press Release 23 April 2023'' \citep[p.~4 (PDF p.~5), n.~7]{cme-24-271}. CME's press release on the conversion is dated 24 April 2023, a Monday: CME ``today announced that it successfully converted 7.5 million contracts of Eurodollar futures and options open interest and \$4 trillion in cleared USD LIBOR swaps to corresponding SOFR derivatives in April'', and only the May and June 2023 contracts remained; it gives no conversion date \citep[reprint dated 24/04/2023]{cme-pr-2023-04-24}.

Three dates are therefore distinct: the \emph{processing weekend} (positions taken after the close on Friday 14 April, transformation on Sunday 16 April), the \emph{trade date} on which the converted positions first existed (Monday 17 April 2023) and the \emph{announcement} (24 April 2023). The ``April 23, 2023'' in the CBOT filing and in CME's 2024 footnote, a Sunday, matches none of them; the filings do not explain it. The press release is dated 24 April.

\qffield{Pitfalls} Eurodollar futures on 3-month LIBOR no longer trade, and EONIA no longer exists. Gliner calls Eurodollars ``the most actively traded futures market in the world based on contact [sic] volume'' (p.~254), as of about 2013, unverified against exchange statistics. \emph{The conversion date}. CME Clearing Advisory 23-082 states the close of business on Friday 14 April 2023 and trade date Monday 17 April 2023. Do not read the filings' ``April 23, 2023'' as either the conversion date or the announcement date. Gliner's fed funds example (FFH4 at 99.79 ``pricing in 0.21'', computed as $100-99.97$) has a typo in the arithmetic (p.~257).

\qffield{See also} \hyperref[def:money-market-reference-rates]{money-market reference rates}; \hyperref[mod:forward-rates-fra]{forward rates and FRAs}; \hyperref[def:interest-rate-swaps]{interest-rate swaps}; \hyperref[def:interest-rate-measures]{interest-rate measures}.

\qffield{Further reading} Market-based policy expectations \citep{rt-kuttner-2001}; \citep{gap-rates-sofr-futures-hedging} on pricing and hedging with SOFR futures.

\end{qfmethod}
